\begin{lemma}[Step 1, Subcase 2.1: two recolorings and the exact prefix]
Let $G$ be a finite loopless multigraph retaining individual parallel-edge
identities, let $e_0$ be distinguished, and let $[k]$ be the color set.
Let $T$ be a positive-rung certified ETT with generating coloring
$\varphi$, and suppose that $\varphi$ is proper and leaves exactly $e_0$
uncolored.  Let $H$ be a good hierarchy over a hierarchy base.  Write
$T_n$ for its last ladder tree, $D_n$ for the computed terminal stored set,
$d$ for the divider count, and $A=H.\mathit{starred}(d)$.  Assume that
$V(T)$ is elementary and that $T$ is
\[
 \left(\bigcup_{\eta\in D^H_{n,d+1}}\Gamma^d_\eta\right)^-
 \text{-closed under }\varphi.
\]

Let $\alpha\ne\beta$, with $\alpha$ missing and singleton-closed on $T$
and $\beta$ not missing on $T$.  Let $P_1,P_2,P_3$ be pairwise
support-disjoint exact $(T,\varphi,\{\alpha,\beta\})$-exit paths.  Their
outside ends $a_1,a_2,a_3$ and tree ends $b_1,b_2,b_3$ are pairwise
distinct, with $b_1\prec_Tb_2\prec_Tb_3$.  For each $j$, retain the exact
$\beta$-colored boundary-edge identity $f_j$ on $P_j$ incident with $b_j$.
Suppose $b_3\notin V(A)$.

Then there are a color $\gamma\ne\beta$, a proper coloring $\mu$ leaving
exactly $e_0$ uncolored, exact prefixes $A\preccurlyeq B^-\prec
C\preccurlyeq T$, and three exact
$(T,\mu,\{\gamma,\beta\})$-exit paths $R_1,R_2,R_3$ with the same
respective endpoints, supports, and lengths as $P_1,P_2,P_3$, such that:
\begin{enumerate}
\item $\mu$ is $(T_n,D_n,\varphi)$-stable, and every exact edge incident
with $B^-$ and every vertex missing-color set on $B^-$ is unchanged;
\item $\gamma$ is absent from the last indexed Gamma union of $H$,
$\gamma$ is missing at $b_3$ under $\varphi$, while $\beta$ is missing
there under $\mu$, and the exact identity $f_3$ has color $\gamma$ under
$\mu$;
\item on every vertex of $C$, membership of every color other than
$\gamma,\beta$ in the missing set is unchanged;
\item the transformed paths remain pairwise support-disjoint;
\item $C$ ends at $b_3$ and is obtained from $B^-$ by one exact
\noderef{TreeSequenceAdjoinEdge} step; and
\item $f_1,f_2$ are still individual $\beta$-colored boundary identities
of $C$, their outside endpoints are distinct, and $b_1,b_2,b_3$ all
belong to $V(C)$.
\end{enumerate}
All prefix relations are exact ordered
\noderef{TreeSequenceInitialSegment} relations.
\end{lemma}
\begin{proof}
Write $D^+=D^H_{n,d+1}$ and
$C_\Gamma=\bigcup_{\eta\in D^+}\Gamma^d_\eta$.  We first locate the
relevant exact prefixes and choose $\gamma$.

The lower and upper fields of \noderef{ExtendedTashkinovTree}, the
nestedness fields of \noderef{TashkinovHierarchy}, and
\noderef{HierarchyStarredStage} give exact ordered-prefix relations
\[
                         T_n\preccurlyeq A\preccurlyeq T.       \tag{1}
\]
For the first relation, if the terminal kind is not PE then the
nonparallel field of \noderef{HierarchyBase} identifies the core with
$T_n$.  In PE, its parallel field and \noderef{OrderedTreeUnion} make
$T_n$ the left ordered segment of the core.  The hierarchy then places
the core below $A$.  Unfolding \noderef{TreeSequenceInitialSegment},
the entrywise vertex and edge equalities compose, so these exact
ordered-prefix relations are transitive.  Applying
\noderef{TreeSequenceVertexSet} to the resulting prefixes gives the
corresponding inclusions of vertex sets.

The initial member of the certified \noderef{TashkinovSeries} is the
initial constructor of \noderef{TashkinovSeriesStep}; its
\noderef{KTriple} field gives $k\geq\Delta(G)+1$.  Under the proper
coloring, \noderef{IsProperEdgeColoring} says that two distinct exact
edges incident with $b_3$, which are adjacent by
\noderef{AdjacentEdges}, cannot have the same nonempty color.
Unfolding \noderef{Incident}, choosing an incident edge witnessing each
member of \noderef{PresentColors} at $b_3$ consequently gives an
injection into the incident-edge set counted by \noderef{Degree}.  By
\noderef{MaximumDegree}, at most $\Delta(G)$ colors are present.  Since
\noderef{MissingColors} is the complementary subset of $[k]$, choose
\[
                         \gamma\in\overline\varphi(b_3).       \tag{2}
\]

From (1) and \noderef{TreeMissingColors},
$\overline\varphi(A)\subseteq\overline\varphi(T)$.  Expanding
\noderef{TashkinovStoredColors}, \noderef{HierarchyStoredColorsAt}, and
\noderef{HierarchyTerminalStoredColors} gives
\[
 D^+=D_n\setminus\overline\varphi(T)
       \subseteq D_n\setminus\overline\varphi(A)=D^H_{n,d}.  \tag{3}
\]
If $\gamma$ belonged to $\Gamma^d_\eta$ for an
$\eta\in D^H_{n,d}$, \noderef{GoodHierarchyReservedDisjoint} would
supply a vertex of $A$ missing $\gamma$.  This vertex differs from
$b_3$, because $b_3\notin V(A)$, while both lie in $V(T)$.  Together
with (2), this contradicts the assumed
\noderef{ColoredVertexSetElementary} property.  Thus $\gamma$ is in no
last indexed Gamma pair, in particular $\gamma\notin C_\Gamma$ by (3).
The terminal \noderef{TreeColorSetMinusClosed} certificate and (2) now
show that $\gamma$ has no colored boundary identity at $T$; the same is
true of $\alpha$ by the assumed \noderef{TreeColorSetClosed} property.
Finally $\gamma\ne\beta$, because $f_3$ is a $\beta$-colored edge
incident with $b_3$, whereas (2), unfolded through
\noderef{PresentColors} and \noderef{MissingColors}, excludes an
incident $\gamma$-edge.

We perform the first recoloring, including its identity branch.  If
$\alpha=\gamma$, put $\mu_4=\varphi$.  Otherwise select every full
$(\alpha,\gamma)$ component containing an $\alpha$-edge node of one of
$P_1,P_2,P_3$, and also the component through $a_j$ when the endpoint
alternative of $P_j$ says that $a_j$ misses $\alpha$.

Each selected component has a vertex-node outside $T$.  Indeed, an
$\alpha$-edge node of a simple exit path is internal and has the two
distinct endpoints of that exact edge as its neighboring vertex-nodes,
by \noderef{TwoColorIncidenceGraph}.  Here adjacency to the edge-node
is exactly endpoint incidence by \noderef{Incident}.  The
unique-tree-vertex clause of \noderef{TreeExitPath} says that either
endpoint in $T$ would have to be $b_j$.  They cannot both be $b_j$,
since the endpoint pair $(b_j,b_j)$ is forbidden by
\noderef{LooplessMultigraph}.  If exactly one were $b_j$, the edge
would be an $\alpha$-colored member of
\noderef{BoundaryEdges}, contrary to singleton-$\alpha$ closedness.
Thus both endpoints are outside $T$.  A component selected through
$a_j$ has the outside vertex $a_j$ itself.

No selected component meets $T$.  Otherwise start from the outside
vertex just found and take a shortest component walk to a vertex-node
of $T$.  The edge-node immediately before the first such vertex has one
end outside and one inside $T$ by the endpoint-incidence definition
\noderef{Incident}, and its color is $\alpha$ or $\gamma$.  It is
therefore an exact boundary identity of one of the two singleton closed
colors, a contradiction.  Finiteness gives an
\noderef{OutsideKempeRecoloring} on precisely these components; call its
result $\mu_4$.

In the unequal branch, \noderef{OutsideKempeRecoloringStable} shows that
$\mu_4$ is proper, leaves exactly $e_0$ uncolored, fixes every exact edge
incident with $T$, preserves every missing set on $T$, and is stable on
$T$ for every control set.  In the equality branch these conclusions
hold by reflexivity, using \noderef{sc2} for stability and the assumed
instance of \noderef{IsOneEdgeUncolored} for the assertion about $e_0$.

For $j=1,2,3$, retain the ordered node list of $P_j$ and call it $Q_j$.
Every old $\beta$-edge stays $\beta$, while every old $\alpha$-edge on
$P_j$ belongs to a selected outside component and becomes $\gamma$.
Thus \noderef{TwoColorIncidenceGraph} gives the required consecutive
adjacencies; simplicity, positive length, endpoints, support, and the
unique vertex-node in $T$ are unchanged.  If $a_j$ missed $\beta$, it
still does, since no edge acquired color $\beta$.  If it missed
$\alpha$, then either $\alpha=\gamma$, or its full component was selected
and every incident edge of either switched color belongs to that
component.  After the swap it therefore misses $\gamma$.  These checks
give
\[
 Q_j:\operatorname{TreeExitPath}(T,\mu_4,\gamma,\beta,a_j,b_j)
 \quad(j=1,2,3),                                               \tag{4}
\]
with the old pairwise-disjoint supports.

Let $K_3$ be the full $(\gamma,\beta)$ component through $b_3$ under
$\mu_4$.  Its support equals the support of $Q_3$.  At $b_3$, (2) and
preservation of its missing set exclude an incident $\gamma$-edge,
while $f_3$ is an incident $\beta$-edge.  The conflict clause of
\noderef{IsProperEdgeColoring}, applied through \noderef{Incident},
excludes a second incident edge of that color, so $b_3$ has degree one
in \noderef{TwoColorIncidenceGraph}.  At $a_3$, (4) makes one color
missing and the positive walk supplies an edge of the other color; the
same properness clause again gives degree one.  Each internal node of
the simple walk has its two distinct walk neighbors, while
\noderef{TwoColorIncidenceDegreeAtMostTwo} excludes a third.  If a node
of $K_3$ lay outside $Q_3$, a shortest component walk from it to the
support of $Q_3$ could attach neither at an internal degree-two node nor
at either degree-one endpoint, whose unique neighbor is already its path
neighbor.  This contradicts connectedness in
\noderef{TwoColorComponent}.  Hence
\[
                         K_3=\operatorname{supp}(Q_3),         \tag{5}
\]
in agreement with \noderef{TwoColorComponentSpanningPath}.  Notice that
we have proved this equality locally; we have not inferred any component
identity through the first recoloring.

Unfold \noderef{TreeVertexOrder}.  Let $s$ be the displayed position of
$b_3$ in $T$.  The two strict order hypotheses put $b_1,b_2$ at earlier
positions.  Writing $a=|A|$, the entrywise vertex equality in
\noderef{TreeSequenceInitialSegment}, together with the vertex
injectivity field of \noderef{TreeSequence}, shows that exact
prefixhood and $b_3\notin V(A)$ give $a<s$.  The positivity field of
\noderef{TreeSequence} and the prefix length inequality give
$1\leq|T_n|\leq a$, so $2\leq s$.  The bounds needed by
\noderef{TreeSequencePrefix} therefore hold for $s-1$ and $s$.  Define
\[
 B^-=\operatorname{TreeSequencePrefix}(T,s-1),\qquad
 C=\operatorname{TreeSequencePrefix}(T,s).                    \tag{6}
\]
The defining entrywise equalities of
\noderef{TreeSequenceInitialSegment}, composed as above, yield
\[
              T_n\preccurlyeq A\preccurlyeq B^-\prec C
                 \preccurlyeq T.                              \tag{7}
\]
The displayed edge $g$ at position $s$ gives exactly
$C=B^-+g$ under \noderef{TreeSequenceAdjoinEdge}; the attachment field
of \noderef{TreeSequence} records its old endpoint, while the prefix
vertex equality and the choice of $s$ show that the new last vertex of
$C$ is $b_3$.  By (5), $K_3$ has no vertex-node in $B^-$.  If an
edge-node of $K_3$ were incident with $B^-$, unfolding
\noderef{TreeIncidentEdge} and \noderef{Incident} would give an
adjacent vertex-node of $K_3$ in $B^-$, a contradiction.  Thus $K_3$
has no edge-node incident with $B^-$; in particular $g\notin K_3$.

Swap $\gamma$ and $\beta$ on $K_3$ and call the result $\mu$.
The inequality $\gamma\ne\beta$ permits
\noderef{TwoColorComponentSwapProper};
\noderef{TwoColorComponentSwap} changes only those two nonempty colors,
so the two clauses of \noderef{IsOneEdgeUncolored} show that exactly
$e_0$ remains uncolored.  Since $K_3$ has no edge-node incident with
$B^-$, every exact edge satisfying \noderef{TreeIncidentEdge} for
$B^-$ keeps its color.  Consequently every \noderef{PresentColors}
set, hence every \noderef{MissingColors} set, at a vertex of $B^-$ also
stays fixed.

These facts prove both clauses of stability on $T_n$.  An edge incident
with $T_n$ has a vertex witness which the entrywise prefix equality in
\noderef{TreeSequenceInitialSegment} transports into $B^-$; hence the
definition \noderef{TreeIncidentEdge} makes it incident with $B^-$, so
its color is fixed whenever the edge clause of
\noderef{StableColoring} asks for it.  Every vertex of $T_n$ similarly
lies in $B^-$, so the vertex clause follows from the missing-set
equality.  Thus $\mu$ is $(T_n,D_n,\mu_4)$-stable.  The corresponding
restriction of the first recoloring is
$(T_n,D_n,\varphi)$-stable; transitivity in \noderef{sc2} gives
\[
             \mu\text{ is }(T_n,D_n,\varphi)\text{-stable}.   \tag{8}
\]

At $b_3$, the unique incident edge-node of $K_3$ is $f_3$.  It changes
from $\beta$ to $\gamma$, and no incident edge changes to $\beta$.
Thus $\mu(f_3)=\gamma$ and
\[
                         \beta\in\overline\mu(b_3).           \tag{9}
\]
For $j=1,2$, support disjointness and (5) put every node of $Q_j$
outside $K_3$, so its edge colors and endpoint alternative are fixed.
For $j=3$, the same simple node list remains an alternating walk because
the two colors exchange throughout its full component.  Its outside-end
missing alternative exchanges in the same way.  Calling the resulting
walks $R_1,R_2,R_3$, we obtain
\[
 R_j:\operatorname{TreeExitPath}(T,\mu,\gamma,\beta,a_j,b_j)
 \quad(j=1,2,3),                                              \tag{10}
\]
with precisely the old supports, lengths, endpoints, and pairwise
support disjointness.

It remains to check the interface facts about $C$.  The one-step
certificate in (6), unfolded through
\noderef{TreeSequenceAdjoinEdge} and
\noderef{TreeSequenceInitialSegment}, says that every vertex of $C$ is
either a vertex of $B^-$ or the new last vertex $b_3$; vertex
injectivity in \noderef{TreeSequence} rules out any further occurrence.
Missing sets agree completely in the first case because every edge
incident with that vertex is a \noderef{TreeIncidentEdge} of $B^-$ and
is fixed.  At $b_3$, the first recoloring fixed every incident edge, and
the second recoloring changed only $\gamma,\beta$.  Membership of every
other color in its missing set is therefore unchanged.  This proves the
asserted pointwise non-$\{\gamma,\beta\}$ comparison on $C$.

The order inequalities place $b_1,b_2$ before the last vertex $b_3$,
so the entrywise prefix relation
\noderef{TreeSequenceInitialSegment} places all three in $V(C)$.  The
original boundary identities $f_1,f_2$ have their inside endpoints
$b_1,b_2$ in $C$, while their other endpoints lie outside $T$ by
\noderef{BoundaryColorEdges} and \noderef{BoundaryEdges}, hence outside
$C$.  The endpoint statements here are incidences in the sense of
\noderef{Incident}; \noderef{LooplessMultigraph} ensures that the two
ends of each boundary edge are distinct.  The first recoloring never
changes a $\beta$-edge.  The second does not change $f_1$ or $f_2$,
because their edge-nodes lie in the supports of $Q_1,Q_2$, which are
disjoint from $K_3$ by (5).  Thus they remain the exact
$\beta$-colored boundary identities of $C$.  Each outside endpoint, by
\noderef{Incident} and the incidence rule of
\noderef{TwoColorIncidenceGraph}, is a vertex-node of the corresponding
retained walk $R_1$ or $R_2$ adjacent to its edge-node $f_1$ or $f_2$;
it is the non-$b_j$ endpoint because it lies outside $C$.  Equality of
those two vertices would
put the same vertex-node in the two supports, contrary to their proved
support disjointness.  Their outside endpoints are therefore distinct.
Equations (2)--(10) and these final observations establish every asserted
clause.
\end{proof}
